% !TEX root = ./2nde.tex
\begin{exo}
	Lors de la Coupe du monde de volley-ball féminin Japon 2018, on donne les tailles des joueuses des équipes des États-Unis et de la Chine.

	\textbf{États-Unis} (en pieds-pouces) :
	\[
		\begin{array}{cccccc}
			5'11'' & 5'9'' & 5'7'' & 6'1'' & 5'10'' & 6'2'' \\
			6'2''  & 6'2'' & 6'4'' & 6'2'' & 6'0''  & 6'4''
		\end{array}
	\]

	\textbf{Chine} (en m) :
	\[
		\begin{array}{ccccccc}
			1,9  & 1,81 & 1,92 & 1,89 & 1,9  & 1,81 & 1,67 \\
			1,82 & 1,72 & 1,87 & 1,96 & 1,93 & 1,82 & 1,84 \\
			1,89 & 1,8  & 1,91 & 1,9  &      &      &
		\end{array}
	\]

	\hfill{\footnotesize\textit{(sources : USA Volley-Ball et Wikipedia.)}}

	\begin{enumerate}
		\item Calculer la taille moyenne et l'écart-type de l'équipe chinoise en mètres.
		\item On sait que 1' = 1 pied = 12 pouces = 12''.
		      \begin{enumerate}
			      \item Convertir chaque taille de l'équipe des États-Unis en pouces.
			      \item Calculer la taille moyenne et l'écart-type de l'équipe états-unienne en pouces.
		      \end{enumerate}
		\item Sachant que 1'' = 2,54 cm, comparer les tailles des deux équipes.
	\end{enumerate}
\end{exo}
